\documentclass[a4paper,11pt]{article}

% --------------------------------
% Sprache, Schrift, Layout
% --------------------------------
\usepackage[T1]{fontenc}
\usepackage[utf8]{inputenc}
\usepackage[german]{babel}
\usepackage{lmodern}
\usepackage{microtype}
\usepackage[a4paper,margin=2.5cm]{geometry}
\usepackage{setspace}
\setstretch{1.15}

% --------------------------------
% Mathematik, Bilder, Links
% --------------------------------
\usepackage{amsmath,amssymb}
\usepackage{graphicx}
\usepackage{xcolor}
\usepackage{float}
\usepackage{hyperref}
\usepackage{caption}

% --------------------------------
% Link-Setup
% --------------------------------
\hypersetup{
    colorlinks=true,
    linkcolor=black,
    urlcolor=blue,
    pdftitle={Übung 1 - Git Basics},
    pdfauthor={Langsenlehner Cornelia}
}

% --------------------------------
% Bildpfade
% --------------------------------
\graphicspath{{screenshots/Ex1/}}

% --------------------------------
% Absatzlayout
% --------------------------------
\setlength{\parindent}{0pt}
\setlength{\parskip}{0.6em}

% --------------------------------
% Eigene Titelseite
% --------------------------------
\begin{document}

\begin{titlepage}
    \centering
    \vspace*{3cm}

    {\scshape\Large Dokumentation der Übung\par}
    \vspace{1cm}

    {\huge\bfseries Git Basics -- PRs, Interactive Rebase \& Unit Tests\par}
    \vspace{2cm}

    {\Large Übung 1\par}
    \vspace{2cm}

    \begin{tabular}{rl}
        \textbf{Name:} & Langsenlehner Cornelia, BSc \\
        \textbf{Matrikelnummer:} & s2510455006 \\
        \textbf{Datum:} & \today \\
    \end{tabular}

    \vfill
\end{titlepage}

\tableofcontents
\newpage

\section{Task 1}

Das Repository wurde erfolgreich geforkt.

\textbf{Forked Repository:}  
\href{https://github.com/CorneliaLan/cd-mcm-exercise-Langsenlehner.git}{github.com/CorneliaLan/cd-mcm-exercise-Langsenlehner.git}

\begin{figure}[H]
    \centering
    \includegraphics[width=0.8\linewidth]{/Users/cornelia/Devs/CI-CD-MCM-Langsenlehner/screenshots/Ex1/fork.png}
    \caption{Geforktes Repository auf GitHub}
    \label{fig:fork}
\end{figure}

Zusätzlich wurde das Projekt ausgeführt und überprüft.

\begin{figure}[H]
    \centering
    \includegraphics[width=0.8\linewidth]{/Users/cornelia/Devs/CI-CD-MCM-Langsenlehner/screenshots/Ex1/run.png}
    \caption{Erfolgreicher Projektlauf}
    \label{fig:run}
\end{figure}

\section{Task 2}

Beschreibung und Dokumentation von Task 2.

\section{Task 3}

Beschreibung und Dokumentation von Task 3.

\section{Task 4}

Beschreibung und Dokumentation von Task 4.

\end{document}

